%! Solving Exponential Equations with a Common Base — Unit 6, Lesson 6.3 — Slides (Beamer)
\documentclass[aspectratio=169, 11pt]{beamer}
\usepackage{algebra2-beamer}   % provides \forestheader{} and \sectionlabel[color]{}
% Standard coordinate grid + axes: {xmin}{xmax}{ymin}{ymax}
\newcommand{\lgrid}[4]{%
  \draw[step=1, linegray!45, thin] (#1,#3) grid (#2,#4);
  \draw[->, thick, charcoal] (#1,0)--(#2,0) node[below right, font=\tiny]{$x$};
  \draw[->, thick, charcoal] (0,#3)--(0,#4) node[left, font=\tiny]{$y$};}
% y = a*b^(x-h) + k plotted as a*exp((ln b)(x-h))+k: pgfmath's pow() is unreliable for
% negative and fractional exponents, but exp() is not.
% #1=a  #2=ln(b)  #3=h  #4=k  #5=xmin  #6=xmax
% ln2 = 0.693147   ln3 = 1.098612
\newcommand{\expgraph}[6]{\draw[very thick, forest, domain=#5:#6, samples=140, smooth]
  plot (\x,{(#1)*exp((#2)*((\x)-(#3)))+(#4)});}
\newcommand{\solline}[3]{\draw[goldacc, very thick] (#1,#3)--(#2,#3);}

\begin{document}

% TITLE SLIDE (CourseName is not defined in beamer, so the name is literal)
{
\setbeamercolor{background canvas}{bg=forest}
\begin{frame}[plain]
  \vspace{2.8cm}\hspace{0.55cm}%
  \begin{minipage}{0.9\textwidth}
    {\color{goldacc}\bfseries\small LESSON 6.3}\\[0.5cm]
    {\color{white}\bfseries\LARGE Solving Exponential Equations with a Common Base}\\[1.2cm]
    {\color{forestmid}\small Algebra 2~$\cdot$~Unit 6: Exponential Functions}
  \end{minipage}
\end{frame}
}

% HOOK
\begin{frame}[t, plain]
  \forestheader{Same height, same question --- two answers or one?}
  \sectionlabel{Hook}
  \begin{columns}[T]
    \begin{column}{0.56\textwidth}
      \centering
      \begin{tikzpicture}[scale=0.26]
        \begin{scope}
          \lgrid{-3}{3}{-1}{9}
          \begin{scope}\clip (-3,-1) rectangle (3,9);
            \draw[very thick, forest, domain=-3:3, samples=90, smooth] plot (\x,{(\x)*(\x)});
          \end{scope}
          \solline{-3}{3}{4}
          \fill[charcoal] (-2,4) circle (4pt);
          \fill[charcoal] (2,4) circle (4pt);
          \node[charcoal, font=\small] at (0,-2.8) {$y=x^{2}$};
        \end{scope}
        \begin{scope}[xshift=8.5cm]
          \lgrid{-3}{3}{-1}{9}
          \begin{scope}\clip (-3,-1) rectangle (3,9); \expgraph{1}{0.693147}{0}{0}{-3}{3} \end{scope}
          \solline{-3}{3}{4}
          \fill[charcoal] (2,4) circle (4pt);
          \node[charcoal, font=\small] at (0,-2.8) {$y=2^{x}$};
        \end{scope}
      \end{tikzpicture}
    \end{column}
    \begin{column}{0.42\textwidth}
      \footnotesize
      \[ x^{2}=4 \qquad\qquad 2^{x}=4 \]
      \begin{itemize}\setlength{\itemsep}{4pt}
        \item The gold line is at the same height in both pictures.
        \item Why does it hit one curve twice and the other once?
      \end{itemize}
      \vspace{3pt}
      \begin{block}{Back to Lesson 6.0}
        An exponential graph has \textbf{no turning points}. It never comes back to an output it has
        already used.
      \end{block}
    \end{column}
  \end{columns}
\end{frame}

% THE PROPERTY
\begin{frame}[t, plain]
  \forestheader{One repeated output would ruin everything --- and there are none}
  \sectionlabel[goldacc]{The permission slip}
  \begin{center}
    \Large If \; $b^{\,m}=b^{\,n}$ \; and \; $b>0$, $b\neq1$, \qquad then \; $m=n$.
  \end{center}
  \vspace{10pt}
  \begin{columns}[T]
    \begin{column}{0.48\textwidth}
      \footnotesize
      \textbf{Why it is true.} An exponential is always increasing or always decreasing --- never both.
      A function that never turns around can never reuse an output, so a horizontal line meets it
      \textbf{at most once}.
    \end{column}
    \begin{column}{0.48\textwidth}
      \begin{block}{What it buys you}
        Once both sides are powers of the \emph{same} base, the bases can be dropped and the exponents
        compared. That is the entire method.
      \end{block}
    \end{column}
  \end{columns}
  \vspace{8pt}
  \begin{center}
    {\footnotesize This is Lesson 6.0's ``no turning points'' doing work, one full unit later --- and
    Unit 7 will use it again.}
  \end{center}
\end{frame}

% THE FINE PRINT
\begin{frame}[t, plain]
  \forestheader{Why $b\neq1$ was never an arbitrary rule}
  \sectionlabel{Fine print}
  \begin{center}
    \Large $1^{5} = 1 \qquad\text{and}\qquad 1^{9} = 1$ \\[6pt]
    \normalsize so $1^{5}=1^{9}$ \; even though \; $5\neq9$.
  \end{center}
  \vspace{10pt}
  \begin{columns}[T]
    \begin{column}{0.48\textwidth}
      \begin{block}{$b=1$ is barred}
        A flat line repeats \emph{every} output. Matching bases would prove nothing about the
        exponents.
      \end{block}
    \end{column}
    \begin{column}{0.48\textwidth}
      \begin{block}{$b>0$ is barred too}
        $(-2)^{1/2}=\sqrt{-2}$ is not real, so a negative base does not even define a curve at every
        exponent.
      \end{block}
    \end{column}
  \end{columns}
  \vspace{8pt}
  \begin{center}
    {\footnotesize In Lesson 6.1 these were rules about \emph{models}. They turn out to be exactly the
    conditions that make today's \emph{solving} legal.}
  \end{center}
\end{frame}

% STEP 1 IS UNIT 5
\begin{frame}[t, plain]
  \forestheader{Step 1 is where the work is --- and it is all Unit 5}
  \sectionlabel{Rewriting}
  \begin{center}
  \footnotesize
  \renewcommand{\arraystretch}{1.35}
  \begin{tabular}{|l|c|c|c|c|c|c|}
  \hline
  \textbf{Powers of $2$} & $4=2^{2}$ & $8=2^{3}$ & $16=2^{4}$ & $32=2^{5}$ & $\frac18=2^{-3}$ & $\sqrt2=2^{1/2}$ \\
  \hline
  \textbf{Powers of $3$} & $9=3^{2}$ & $27=3^{3}$ & $81=3^{4}$ & $\frac13=3^{-1}$ & $\frac1{27}=3^{-3}$ & $1=3^{0}$ \\
  \hline
  \end{tabular}
  \end{center}
  \vspace{8pt}
  \begin{columns}[T]
    \begin{column}{0.48\textwidth}
      \begin{block}{Multiply, don't add}
        $\left(2^{3}\right)^{x}=2^{\,3x}$ \\[3pt]
        {\scriptsize Check: $\left(2^{3}\right)^{4}=4096=2^{12}$, while $2^{3}\cdot2^{4}=128=2^{7}$.}
      \end{block}
    \end{column}
    \begin{column}{0.48\textwidth}
      \begin{block}{Reach the whole exponent}
        $\left(3^{2}\right)^{x+4}=3^{\,2(x+4)}=3^{\,2x+8}$ \\[3pt]
        {\scriptsize Writing $3^{\,2x+4}$ multiplies the $x$ and forgets the $4$.}
      \end{block}
    \end{column}
  \end{columns}
\end{frame}

% THE METHOD
\begin{frame}[t, plain]
  \forestheader{Three steps, twice}
  \sectionlabel[goldacc]{The method}
  \begin{center}
    {\footnotesize \textbf{1.} Rewrite both sides over one base. \quad \textbf{2.} Set the exponents
    equal. \quad \textbf{3.} Solve and check.}
  \end{center}
  \vspace{6pt}
  \begin{columns}[T]
    \begin{column}{0.48\textwidth}
      \footnotesize
      \textbf{A.} \; $2^{\,x+1}=32$
      \begin{align*}
        2^{\,x+1} &= 2^{5} \\
        x+1 &= 5 \\
        x &= 4
      \end{align*}
    \end{column}
    \begin{column}{0.48\textwidth}
      \footnotesize
      \textbf{B.} \; $8^{x}=4$
      \begin{align*}
        \left(2^{3}\right)^{x} = 2^{\,3x} &= 2^{2} \\
        3x &= 2 \\
        x &= \tfrac23
      \end{align*}
    \end{column}
  \end{columns}
  \vspace{4pt}
  \begin{center}
    \begin{minipage}{0.86\textwidth}
      \begin{block}{A fraction is not a mistake}
        $8^{\,2/3}=\left(\sqrt[3]{8}\right)^{2}=2^{2}=4$ \checkmark \\[3pt]
        {\scriptsize The exponent is a location on the $x$-axis. Nothing requires it to be a whole
        number --- and students who ``fix'' it are erasing correct work.}
      \end{block}
    \end{minipage}
  \end{center}
\end{frame}

% HARDER
\begin{frame}[t, plain]
  \forestheader{When both sides have to move}
  \sectionlabel{Harder}
  \begin{columns}[T]
    \begin{column}{0.46\textwidth}
      \footnotesize
      \textbf{C.} \; $3^{x}=\dfrac{1}{27}$
      \begin{align*}
        3^{x} &= 3^{-3} \\
        x &= -3
      \end{align*}
      \vspace{2pt}
      \textbf{D.} \; $9^{\,x-2}=27^{x}$
      \begin{align*}
        3^{\,2(x-2)} &= 3^{\,3x} \\
        2x-4 &= 3x \\
        x &= -4
      \end{align*}
    \end{column}
    \begin{column}{0.50\textwidth}
      \begin{block}{The two errors of the day}
        \footnotesize
        $\left(2^{3}\right)^{x}=2^{\,3+x}$ \; --- added instead of multiplied. \\[4pt]
        $9^{\,x-2}=3^{\,2x-2}$ \; --- distributed to the $x$ only. \\[6pt]
        Both come out as clean-looking wrong answers. Write the parentheses \emph{before} removing
        them.
      \end{block}
      \vspace{4pt}
      {\footnotesize Check D: $9^{-6}=3^{-12}$ and $27^{-4}=3^{-12}$. \checkmark}
    \end{column}
  \end{columns}
\end{frame}

% INTERSECTIONS
\begin{frame}[t, plain]
  \forestheader{Every solution is a crossing point}
  \sectionlabel{Graphically}
  \begin{columns}[T]
    \begin{column}{0.38\textwidth}
      \centering
      \begin{tikzpicture}[scale=0.30]
        \lgrid{-3}{4}{-4}{9}
        \begin{scope}\clip (-3,-4) rectangle (4,9); \expgraph{1}{0.693147}{0}{0}{-3}{4} \end{scope}
        \solline{-3}{4}{4}
        \draw[redacc, dashed, very thick] (-3,-3)--(4,-3);
        \node[goldacc, font=\tiny, right] at (2.3,4.8) {$y=4$};
        \node[redacc, font=\tiny, right] at (2.3,-3.7) {$y=-3$};
        \fill[charcoal] (2,4) circle (5pt);
      \end{tikzpicture}
    \end{column}
    \begin{column}{0.58\textwidth}
      \footnotesize
      \begin{itemize}\setlength{\itemsep}{4pt}
        \item $2^{x}=4$ \; is answered by the crossing at $x=2$ --- the same point where $y=2^{x}-4$
              met the $x$-axis in Lesson 6.2.
        \item $2^{x}=-3$ \; has \textbf{no solution}: the range of $y=2^{x}$ is $(0,\infty)$, and the
              curve never reaches, let alone crosses, its asymptote.
      \end{itemize}
      \vspace{3pt}
      \begin{block}{Count the answers}
        For $b^{x}=c$: \textbf{one} solution when $c>0$, \textbf{none} when $c\le0$. Never two --- the
        $\pm$ belongs to the parabola.
      \end{block}
    \end{column}
  \end{columns}
\end{frame}

% THE WALL
\begin{frame}[t, plain]
  \forestheader{$2^{x}=5$: the line crosses, and we still cannot say where}
  \sectionlabel[goldacc]{The wall}
  \begin{columns}[T]
    \begin{column}{0.38\textwidth}
      \centering
      \begin{tikzpicture}[scale=0.30]
        \lgrid{-2}{4}{-1}{10}
        \begin{scope}\clip (-2,-1) rectangle (4,10); \expgraph{1}{0.693147}{0}{0}{-2}{4} \end{scope}
        \solline{-2}{4}{5}
        \node[goldacc, font=\tiny, left] at (-0.4,5.7) {$y=5$};
        \fill[charcoal] (2,4) circle (4pt) node[right, font=\tiny]{$(2,4)$};
        \fill[charcoal] (3,8) circle (4pt) node[right, font=\tiny]{$(3,8)$};
      \end{tikzpicture}
    \end{column}
    \begin{column}{0.58\textwidth}
      \footnotesize
      The solution \textbf{exists} --- it sits between $x=2$ and $x=3$. But $5$ is not a power of $2$,
      so Step 1 is impossible.
      \vspace{4pt}
      \begin{block}{Two different failures, kept apart}
        \textbf{No solution:} $2^{x}=-3$. The graphs never meet. \\[3pt]
        \textbf{Out of reach:} $2^{x}=5$. The graphs meet somewhere our algebra cannot name.
      \end{block}
      \vspace{4pt}
      {\scriptsize Add to the list: the coffee's $132(0.85)^{t}=32$ and the soda's $(0.9)^{t}=2.118$
      from Lesson 6.2. All three have answers. None has a common base.}
    \end{column}
  \end{columns}
\end{frame}

% ROADMAP
\begin{frame}[t, plain]
  \forestheader{Where Unit 6 goes next}
  \sectionlabel{Connections}
  \textbf{Today:} rewrite both sides over one base, set the exponents equal, solve --- legal because an
  exponential never repeats an output. A solution is where a horizontal line crosses the curve, and a
  line below the $x$-axis never crosses at all.\\[8pt]
  \begin{itemize}\setlength{\itemsep}{3pt}
    \item \textbf{6.4} Compound interest, half-life, and $e$ --- real models, and $1000(1.05)^{t}=2000$, where $2$ is not a power of $1.05$.
    \item \textbf{6.5} Exponential regression --- technology fits the model; you decide whether it is the right family.
    \item \textbf{Unit 7} Logarithms --- the tool that finally names the crossing point of $y=2^{x}$ and $y=5$.
  \end{itemize}
  \vspace{6pt}
  \begin{center}
    {\footnotesize The method taught today is narrow on purpose. Knowing exactly \emph{where} it stops
    is what makes the next two lessons make sense.}
  \end{center}
\end{frame}

\end{document}
